%% Beispiel-Lebenslauf - Ana Lorena Ortiz Loyola
%% Umfassende Übersicht aller Kompetenzen, Erfahrungen und Erfolge.
%% Als Master-Referenz für die Erstellung rollenspezifischer Lebensläufe verwenden.
%% Deutsche Version von main_example.tex.
%%
%% Kompilieren mit: lualatex main_example_de.tex
%% (pdflatex schlägt auf modernem MiKTeX häufig mit fontawesome5-Fehlern fehl.)

\documentclass[11pt,a4paper,sans]{moderncv}
\moderncvstyle{banking}
\moderncvcolor{blue}

% Erzwingt, dass Vor- und Nachname sowie die Sektionsüberschriften in moderncv
% Blau (color1) dargestellt werden. Der banking-Standardstil rendert diese auf
% lualatex+MiKTeX sonst schwarz, was inkonsistent mit dem restlichen blauen
% Akzentschema wirkt.
\renewcommand*{\firstnamestyle}[1]{{\fontsize{34}{36}\bfseries\upshape\color{color1}#1}}
\renewcommand*{\lastnamestyle}[1]{{\fontsize{34}{36}\bfseries\upshape\color{color1}#1}}
\renewcommand*{\sectionstyle}[1]{{\sectionfont\color{color1}#1}}

\usepackage[utf8]{inputenc}
\usepackage{hyperref}
\hypersetup{
    colorlinks=true,
    linkcolor=blue,
    filecolor=magenta,
    urlcolor=blue,
    pdftitle={Ana Lorena Ortiz Loyola - Lebenslauf},
    pdfpagemode=FullScreen,
}
\usepackage[scale=0.80]{geometry}
\usepackage{import}
\usepackage{needspace}

% Persönliche Daten
\name{Ana Lorena}{Ortiz Loyola}
\address{Frankfurt am Main, Deutschland}{}{}
\phone[mobile]{+49 176 62145797}
\email{a01750283@tec.mx}
\extrainfo{\href{https://linkedin.com/in/lorenasolana}{LinkedIn}}

\begin{document}

\makecvtitle

% ============================================================
%     PROFILTEXT
% ============================================================

\vspace{6pt}
\small{Data-Science-Studentin (Tecnológico de Monterrey, Abschluss voraussichtlich 2027) mit praktischer Erfahrung in Data Engineering, Machine Learning und kryptografischer Forschung. Entwickelte automatisierte Datenqualitäts- und Migrationstests bei der Deutschen Bundesbank (AnaCredit) sowie eigenständige Projekte in den Bereichen Predictive Modeling, Klassifikation und gitterbasierte Kryptografie. Übernimmt sicher die Verantwortung für eine Datenpipeline end-to-end, von der Datenerfassung bis zur revisionssicheren Dokumentation.}

% ============================================================
%     KERNKOMPETENZEN
% ============================================================

\section{Kernkompetenzen}
\vspace{1pt}
\begin{itemize}

\item \textbf{Programmierung}: Python (Pandas, NumPy, Scikit-Learn, Keras, TensorFlow), R, Java, C++

\item \textbf{Daten \& Datenbanken}: SQL, SQLite, Datenbereinigung und Pipeline-Design, Excel

\item \textbf{Machine Learning}: Random Forest, SVM, K-Means, Predictive Modeling, Klassifikation, Clustering

\item \textbf{Kryptografie \& Sicherheit}: Gitterbasierte Verfahren (LWE, NTRU), Komplexitäts- und Sicherheitsparameteranalyse

\item \textbf{Cloud \& Tools}: Azure Cognitive Services, AWS, Jira, Git, Tableau, Matplotlib, Seaborn

\end{itemize}

% ============================================================
%     BERUFSERFAHRUNG
% ============================================================

\section{Berufserfahrung}
\vspace{3pt}
\begin{itemize}

% --- Aktuellste Position ---
\item{\cventry{04/2026--07/2026}{Data Engineer / Praktikantin}{Deutsche Bundesbank (AnaCredit)}{Frankfurt am Main, Deutschland}{}{\vspace{1pt}
\begin{itemize}
    \item Entwickelte automatisierte Migrationstests (Zeilenanzahl, referenzielle Integrität, Aggregatkonsistenz) mit Pandas-DataFrames, um identische Geschäftslogik nach der Datenmigration sicherzustellen
    \item Implementierte automatisierte Fehlererkennung zur Identifikation von Abweichungen (doppelte IDs, Wertebereichsverletzungen)
    \item Automatisierte die Erstellung von Jira-Tickets sowie E-Mails an meldende Banken bei festgestellten Datenqualitätsproblemen
    \item Reduzierte die Reaktionszeit bei Datenqualitätsproblemen deutlich
    \item Arbeitete im agilen Scrum-Framework (Sprint-Planung, Daily Stand-ups, Sprint Reviews)
\end{itemize}
\vspace{3pt}
\textbf{Automatisiertes Validierungsframework für Meldewesen-Daten} -- überführte einen nicht skalierbaren, fallweisen Validierungstest in eine parametrisierte, wiederverwendbare Python-Lösung
\begin{itemize}
    \item Entwarf parametrisierte MySQL-Abfragen, die dynamisch auf beliebige Datenbanken, Attribute und Meldezeiträume anwendbar sind -- anstelle fest codierter Einzelfall-Skriptlogik
    \item Verlagerte rechenintensive Operationen auf die Datenbankebene (SQL-Views, aggregierte Abfragen) zur effizienten Verarbeitung großer aufsichtlicher Datenbanken
    \item Entwickelte eine YAML-basierte Konfigurationsschnittstelle für Batch-Validierungen über mehrere Datenbanken, Attribute und Meldezeiträume in einem einzigen Aufruf
    \item Erweiterte die Validierungslogik um perzentilbasierte statistische Benchmarks über ein rollierendes 12-Monats-Fenster, um auffällige Meldefrequenzen zu erkennen und echte Datenqualitätsprobleme von falsch-positiven Befunden zu unterscheiden
    \item Implementierte umfassende Fehlerbehandlung für Verbindungsabbrüche, ungültige Eingaben und Grenzfälle wie zusammengesetzte Primärschlüssel
\end{itemize}
\vspace{3pt}
\textbf{Testdatengenerator für die Bundesbank-Plattform} -- entwickelte ein konfigurationsgesteuertes System zur Erzeugung vollständiger, konsistenter synthetischer Testdatensätze
\begin{itemize}
    \item Erzeugte kartesische Attributkombinationen je Tabelle aus einer benutzerdefinierten Vorlage, unabhängig vom Datentyp (Strings, Booleans, Zahlen, Datumsformate)
    \item Hielt IDs tabellenübergreifend konsistent, um referenzielle Integrität in den generierten Daten sicherzustellen
    \item Exportierte Datensätze als XML mit vollständigen Schema-Metadaten, Informationen zur referenziellen Integrität und Plausibilitätskennzahlen (z. B. erwartete Datensatzanzahlen)
\end{itemize}}}

\end{itemize}

% ============================================================
%     PROJEKTE
% ============================================================

\section{Projekte}
\vspace{1pt}
\begin{itemize}

\needspace{5\baselineskip}
\item{\cventry{10/2024--02/2025}{Forschung zu gitterbasierter Kryptografie}{Tecnológico de Monterrey}{}{}{\vspace{1pt}
\begin{itemize}
    \item Implementierte und analysierte die Komplexität gitterbasierter kryptografischer Algorithmen (LWE, NTRU) in Python
    \item Trug zu einer wissenschaftlichen Publikation bei, die Sicherheitsparameter und Laufzeitverhalten über verschiedene Gitterdimensionen hinweg analysiert
\end{itemize}}}

\vspace{3pt}

\item{\cventry{10/2024--12/2024}{Predictive Modeling -- Optimierung im Gesundheitswesen}{Tecnológico de Monterrey}{}{}{\vspace{1pt}
\begin{itemize}
    \item Analysierte über 25 Jahre mexikanischer nationaler Gesundheitsdaten zur Identifikation von Patientenrisikoprofilen und Ressourcenengpässen
    \item Entwickelte Predictive-Modelle (Random Forest, neuronale Netze) und K-Means-Clustering zur Segmentierung von Hochrisikogruppen
\end{itemize}}}

\vspace{3pt}

\item{\cventry{05/2026}{Banking-Datenpipeline -- Datenbereinigung \& Prüfprotokoll}{Persönliches Projekt (GitHub)}{}{}{\vspace{1pt}
\begin{itemize}
    \item Entwickelte eine End-to-End-Pipeline für simulierte Banktransaktionen (150 Konten) mit einer normalisierten SQLite-Datenbank und Fremdschlüssel-Integritätsprüfungen
    \item Mehrstufige Datenbereinigung mit vollständig nachvollziehbarem Protokoll; automatisierte Kontostandspflege über SQL-Trigger
\end{itemize}}}

\vspace{3pt}

\item{\cventry{04/2023--06/2023}{Luftqualitätsanalyse Mexiko-Stadt -- SVM-Klassifikation}{Tecnológico de Monterrey}{}{}{\vspace{1pt}
\begin{itemize}
    \item Analysierte stadtweite NOx/O3-Daten und wandte SVM-Modelle (linear \& RBF) zur Klassifikation kritischer Ozonwerte mit hoher Genauigkeit an
    \item Korrelationsanalysen und Visualisierungen identifizierten erhöhte Belastung in südlichen Stadtbezirken; Ergebnisse in Abschlusspräsentation kommuniziert
\end{itemize}}}

\end{itemize}

% ============================================================
%     AUSBILDUNG
% ============================================================

\section{Ausbildung}
\vspace{1pt}
\begin{itemize}

\item{\cventry{2022--2027}{BSc in Data Science}{Tecnológico de Monterrey}{Mexiko}{}{\vspace{1pt}
Note: 1,3 (Stipendium für akademische Exzellenz; DAAD-KOSPIE-Stipendium). Schwerpunkte: Datenanalyse \& Statistik, Machine Learning \& KI, Kryptografie \& Sicherheit, NLP, numerische Optimierung.
}}

\vspace{3pt}

\item{\cventry{2025--2026}{Auslandssemester}{Universität Göttingen}{Deutschland}{}{\vspace{1pt}
}}

\vspace{3pt}

\item{\cventry{2019--2022}{Abitur}{Tecnológico de Monterrey}{Mexiko}{}{\vspace{1pt}
Note: 1,0. Nationale Physik-Olympiade 2020; Lobende Erwähnung, Physik-Olympiade 2021.
}}

\end{itemize}

% ============================================================
%     SPRACHEN
% ============================================================

\section{Sprachen}
\vspace{1pt}
\begin{itemize}
\item Spanisch (Muttersprache), Englisch (C1, TOEFL iBT 105/120), Deutsch (C1).
\end{itemize}

% ============================================================
%     PUBLIKATIONEN (optional)
% ============================================================

\section{Publikationen}
\vspace{1pt}
\begin{itemize}
\item Beitrag zu einer wissenschaftlichen Publikation über gitterbasierte Kryptografie (Sicherheitsparameter- und Laufzeitanalyse von LWE und NTRU), gemeinsam mit Prof. A. F. De Abiega L'Eglisse, Tecnológico de Monterrey (Zitation ausstehend, Publikation in Vorbereitung).
\end{itemize}

% ============================================================
%     AUSZEICHNUNGEN (optional)
% ============================================================

\section{Auszeichnungen}
\vspace{1pt}
\begin{itemize}
\item{\textbf{Stipendium für akademische Exzellenz} -- Tecnológico de Monterrey.}
\item{\textbf{DAAD-KOSPIE-Stipendium}.}
\item{\textbf{Nationale Physik-Olympiade} -- Mexiko (2020).}
\item{\textbf{Lobende Erwähnung, Physik-Olympiade} (2021).}
\end{itemize}

% ============================================================
%     REFERENZEN
% ============================================================

\section{Referenzen}
\vspace{1pt}
\begin{itemize}
\item{Auf Anfrage erhältlich.}
\end{itemize}

\end{document}
